\documentclass[12pt, a4paper]{article}

% ============================================================
% PACCHETTI
% ============================================================
\usepackage[utf8]{inputenc}
\usepackage[T1]{fontenc}
\usepackage[italian]{babel}
\usepackage{lmodern}
\usepackage{microtype}
\usepackage{geometry}
\geometry{
    top=2.5cm, bottom=2.5cm,
    left=3cm, right=3cm
}

% Grafica e colori
\usepackage{graphicx}
\usepackage[table,xcdraw]{xcolor}
\definecolor{insubriablue}{RGB}{0,80,150}
\definecolor{codeblue}{RGB}{30,100,180}
\definecolor{codegray}{RGB}{245,245,245}
\definecolor{codered}{RGB}{180,30,30}
\definecolor{codegreen}{RGB}{0,120,60}

% Matematica
\usepackage{amsmath, amssymb, amsthm}
\usepackage{bm}

% Tabelle
\usepackage{booktabs}
\usepackage{longtable}
\usepackage{multirow}
\usepackage{array}

% Codice
\usepackage{listings}
\lstdefinestyle{pythonstyle}{
    backgroundcolor=\color{codegray},
    commentstyle=\color{codegreen}\itshape,
    keywordstyle=\color{codeblue}\bfseries,
    stringstyle=\color{codered},
    basicstyle=\ttfamily\small,
    breakatwhitespace=false,
    breaklines=true,
    keepspaces=true,
    numbers=left,
    numberstyle=\tiny\color{gray},
    numbersep=8pt,
    showspaces=false,
    showstringspaces=false,
    tabsize=4,
    frame=single,
    rulecolor=\color{gray!40},
    language=Python
}
\lstset{style=pythonstyle}

% Link e riferimenti
\usepackage[colorlinks=true,
            linkcolor=insubriablue,
            citecolor=insubriablue,
            urlcolor=insubriablue]{hyperref}
\usepackage{cleveref}

% Intestazione e piè di pagina
\usepackage{fancyhdr}
\pagestyle{fancy}
\fancyhf{}
\fancyhead[L]{\small\textcolor{insubriablue}{Pipeline EEG -- Iowa Gambling Task}}
\fancyhead[R]{\small\textcolor{insubriablue}{Interfacce Uomo-Macchina a.a.\ 2025/26}}
\fancyfoot[C]{\thepage}
\renewcommand{\headrulewidth}{0.4pt}

% Paragrafo
\usepackage{parskip}
\setlength{\parskip}{6pt}
\usepackage{enumitem}

% Callout box
\usepackage{tcolorbox}
\tcbuselibrary{skins, breakable}
\newtcolorbox{notabox}[1][]{
    enhanced, breakable,
    colback=insubriablue!5,
    colframe=insubriablue!60,
    fonttitle=\bfseries,
    title=#1
}
\newtcolorbox{formulabox}[1][]{
    enhanced,
    colback=orange!5,
    colframe=orange!60,
    fonttitle=\bfseries,
    title=#1
}

% ============================================================
% DOCUMENTO
% ============================================================
\begin{document}

% ---- FRONTESPIZIO ----
\begin{titlepage}
\centering
\vspace*{1cm}
{\large Università degli Studi dell'Insubria\\
Laurea Triennale in Informatica}\\[0.5cm]
{\large Corso: \textbf{Interfacce Uomo-Macchina} -- a.a.\ 2025/26}\\[0.3cm]
{\small Docente: Prof.ssa Silvia Corchs}\\[2cm]

\rule{\textwidth}{1.5pt}\\[0.4cm]
{\LARGE \textbf{Documentazione Tecnica}}\\[0.3cm]
{\Large Pipeline di Preprocessing EEG}\\[0.2cm]
{\large per la Classificazione di Decisioni\\
nell'Iowa Gambling Task}\\[0.4cm]
\rule{\textwidth}{1.5pt}\\[2cm]

{\large \textbf{Dataset di riferimento:}}\\[0.2cm]
Chávez-Sánchez M., Torres-Ramos S., Román-Godínez I., Salido-Ruiz R.A.\\
\textit{Behavioral and electroencephalographic dataset simultaneously acquired\\
during the Iowa Gambling Task}\\
Scientific Data, 13:359, 2026\\
\href{https://doi.org/10.1038/s41597-026-06662-0}{DOI: 10.1038/s41597-026-06662-0}\\[1.5cm]

{\large \textbf{Dataset disponibile su:}}\\
\href{https://data.mendeley.com/datasets/2pw2m39yct/2}{Mendeley Data -- 2pw2m39yct, versione 2}\\[2cm]

\vfill
{\large Progetto d'esame -- Parte I: Preprocessing EEG}\\[0.3cm]
{\large Maggio 2026}
\end{titlepage}

% ---- INDICE ----
\tableofcontents
\newpage

% ============================================================
\section{Introduzione e Obiettivi}
% ============================================================

\subsection{Contesto del progetto}

Il presente documento descrive nel dettaglio la pipeline di preprocessing EEG implementata nell'ambito del progetto d'esame del corso di Interfacce Uomo-Macchina. L'obiettivo generale del progetto è sviluppare un sistema capace di \textbf{classificare automaticamente le decisioni umane} -- distinguendo decisioni \textit{vantaggiose} da decisioni \textit{svantaggiose} -- a partire dal segnale elettroencefalografico (EEG) acquisito durante l'esecuzione dell'Iowa Gambling Task (IGT).

Il preprocessing costituisce la \textbf{fase fondante} dell'intera pipeline: la qualità del segnale elaborato determina direttamente le performance dei classificatori a valle. Un segnale non correttamente filtrato, non dereferenziato o contaminato da artefatti produce feature spurie che degradano qualsiasi modello di machine learning, indipendentemente dalla sua sofisticazione.

\subsection{Obiettivo della classificazione}

\begin{notabox}[Task di classificazione]
\textbf{Input:} segnale EEG nella finestra temporale $[-2\text{s},\, 0\text{s}]$ prima del momento di decisione.\\[4pt]
\textbf{Output:} classe binaria
\begin{itemize}[noitemsep]
    \item \textbf{Classe 0} -- Decisione \textit{svantaggiosa} (deck A o B)
    \item \textbf{Classe 1} -- Decisione \textit{vantaggiosa} (deck C o D)
\end{itemize}
\end{notabox}

\subsection{Struttura della documentazione}

Il documento è organizzato come segue:
\begin{enumerate}
    \item \textbf{Sezione 2:} descrizione del dataset (paper, acquisizione, struttura file)
    \item \textbf{Sezione 3:} pipeline di preprocessing passo per passo
    \item \textbf{Sezione 4:} sincronizzazione EEG--IGT e generazione label
    \item \textbf{Sezione 5:} estrazione delle epoche
    \item \textbf{Sezione 6:} output e formato dati per ML
    \item \textbf{Sezione 7:} risultati ottenuti e confronto con la letteratura
    \item \textbf{Sezione 8:} scelte implementative e motivazioni
\end{enumerate}

% ============================================================
\section{Dataset}
% ============================================================

\subsection{Descrizione generale}

Il dataset utilizzato è stato pubblicato da Chávez-Sánchez et al.\ (2026) su \textit{Scientific Data} ed è pubblicamente disponibile su Mendeley Data. Contiene registrazioni simultanee di dati comportamentali EEG raccolte durante l'esecuzione dell'Iowa Gambling Task su 59 partecipanti non clinici.

\begin{table}[h!]
\centering
\caption{Caratteristiche principali del dataset}
\begin{tabular}{ll}
\toprule
\textbf{Parametro} & \textbf{Valore} \\
\midrule
Numero partecipanti & 59 (30 M, 29 F) \\
Età media & 21.44 anni ($\sigma = 1.30$) \\
Gruppi & Engineering (ENG, $n=29$), Physical Culture and Sports (PCS, $n=30$) \\
Dispositivo EEG & Natus XLTEK EEG 32® \\
Elettrodi attivi & 21 (sistema internazionale 10-20) \\
Riferimento acquisizione & Pz \\
Elettrodo di terra & CPz \\
Frequenza di campionamento & 256 Hz \\
Impedenza minima & 10 k$\Omega$ \\
Trial IGT per soggetto & 200 (2 blocchi da 100) \\
Durata media sessione & $\approx 16$ minuti \\
\bottomrule
\end{tabular}
\end{table}

\subsection{Protocollo sperimentale}

La sessione sperimentale è strutturata in quattro momenti distinti:

\begin{enumerate}
    \item \textbf{Baseline (3 min):} il partecipante mantiene gli occhi aperti in stato di riposo; il segnale EEG viene registrato come riferimento basale.
    \item \textbf{Primo blocco IGT (100 decisioni):} il partecipante esegue il task scegliendo tra i deck A, B, C, D.
    \item \textbf{Pausa (4 min):} periodo di riposo tra i due blocchi.
    \item \textbf{Secondo blocco IGT (100 decisioni):} completamento del task.
\end{enumerate}

Ogni iterazione del task si compone di due fasi:
\begin{itemize}
    \item \textbf{Pre-decisione:} il partecipante delibera quale deck scegliere (durata variabile, dipendente dal RT individuale).
    \item \textbf{Feedback (1500 ms):} il partecipante visualizza il risultato (crediti vinti/persi).
\end{itemize}

\subsection{Iowa Gambling Task -- Struttura dei payoff}

L'IGT è stato progettato da Bechara et al.\ (1994) per valutare il processo decisionale in condizioni di incertezza. I quattro deck hanno strutture di payoff asimmetriche:

\begin{table}[h!]
\centering
\caption{Struttura dei payoff dei deck nell'IGT originale}
\begin{tabular}{ccccc}
\toprule
\textbf{Deck} & \textbf{Guadagno/carta} & \textbf{Perdita occasionale} & \textbf{Netto ogni 10 carte} & \textbf{Tipo} \\
\midrule
A & +100 & $-\{150,200,250,300,350\}$ & $-250$ & Svantaggioso \\
B & +100 & $-1250$ (1 volta su 10) & $-250$ & Svantaggioso \\
C & +50  & $-\{25,50,75\}$ & $+250$ & Vantaggioso \\
D & +50  & $-250$ (1 volta su 10) & $+250$ & Vantaggioso \\
\bottomrule
\end{tabular}
\end{table}

I deck A e B producono guadagni immediati elevati ma perdite nette a lungo termine ($-250$ ogni 10 carte). I deck C e D producono guadagni più modesti ma un bilancio positivo nel lungo periodo ($+250$ ogni 10 carte). Un soggetto sano e non impulsivo dovrebbe progressivamente preferire C e D.

\subsection{Struttura dei file}

Il dataset è organizzato secondo lo standard BIDS (\textit{Brain Imaging Data Structure}):

\begin{verbatim}
root/
├── participants.csv           (demografici: sesso, età, gruppo)
├── obstetrics.csv             (dati ostetrici per le partecipanti donne)
├── substance_use.csv          (uso di sostanze)
├── dataset-description.csv
├── s-01/
│   ├── EEG.csv               (segnale raw, 21 colonne = canali)
│   ├── IGT.csv               (dati comportamentali)
│   └── processed_EEG.csv     (EEG preprocessato dagli autori)
├── s-02/
│   └── ...
└── s-59/
\end{verbatim}

\subsubsection{Formato EEG.csv}
Matrice di dimensioni $(N_\text{campioni} \times 21)$, dove ogni colonna corrisponde a un elettrodo. I valori sono in microvolt ($\mu$V). Per il soggetto s-01: $N_\text{campioni} = 270336$, corrispondente a 1056 secondi ($\approx 17.6$ minuti) a 256 Hz.

\subsubsection{Formato IGT.csv}
Ogni riga corrisponde a un trial. Le colonne rilevanti sono:

\begin{table}[h!]
\centering
\caption{Colonne del file IGT.csv}
\begin{tabular}{lll}
\toprule
\textbf{Colonna} & \textbf{Tipo} & \textbf{Descrizione} \\
\midrule
\texttt{iteration} & int & Numero del trial (1--200) \\
\texttt{EEG sample} & int & Indice campione EEG sincronizzato con la decisione \\
\texttt{decision} & str & Deck scelto: A, B, C o D \\
\texttt{win} & int & Crediti vinti \\
\texttt{lose} & int & Crediti persi (valore negativo) \\
\texttt{balance} & int & Bilancio cumulativo \\
\bottomrule
\end{tabular}
\end{table}

\subsection{Canali EEG e posizionamento}

I 21 elettrodi attivi, referenziati a Pz durante l'acquisizione, coprono le principali regioni corticali:

\begin{table}[h!]
\centering
\caption{Elettrodi EEG e aree corticali associate}
\begin{tabular}{ll}
\toprule
\textbf{Regione} & \textbf{Elettrodi} \\
\midrule
Frontale & Fp1, Fp2, Fpz, F3, F4, F7, F8, Fz \\
Centrale & C3, C4, Cz \\
Temporale & T4, T5, T6, A1, A2 \\
Parietale & P3, P4, Pz \\
Occipitale & O1, O2 \\
\bottomrule
\end{tabular}
\end{table}

% ============================================================
\section{Pipeline di Preprocessing EEG}
% ============================================================

La pipeline è implementata in Python 3.10+ utilizzando la libreria MNE-Python (versione $\geq$ 1.6) come framework principale per l'elaborazione EEG. La struttura modulare è organizzata in funzioni indipendenti e componibili.

\begin{notabox}[Panoramica della pipeline]
\textbf{Input:} \texttt{EEG.csv} + \texttt{IGT.csv} per ogni soggetto\\[4pt]
\textbf{Step 1:} Caricamento e creazione oggetto MNE Raw\\
\textbf{Step 2:} Re-referencing (Common Average Reference)\\
\textbf{Step 3:} Filtraggio bandpass 0.5--70 Hz\\
\textbf{Step 4:} Filtraggio notch 50 Hz\\
\textbf{Step 5:} Rimozione artefatti via ICA\\
\textbf{Step 6:} Normalizzazione z-score per canale\\
\textbf{Step 7:} Sincronizzazione EEG--IGT\\
\textbf{Step 8:} Estrazione epoche $[-2\text{s},\, 0\text{s}]$\\
\textbf{Step 9:} Generazione label e salvataggio\\[4pt]
\textbf{Output:} tensori \texttt{.npy} pronti per classificazione ML
\end{notabox}

% ----
\subsection{Step 1 -- Caricamento e creazione oggetto MNE Raw}

\subsubsection{Caricamento}
Il file \texttt{EEG.csv} viene letto con \texttt{pandas.read\_csv()} producendo una matrice NumPy di shape $(N_\text{campioni},\, 21)$. I valori sono in $\mu$V come forniti dal dispositivo di acquisizione.

\subsubsection{Conversione in MNE RawArray}
MNE-Python richiede che i dati siano in \textbf{Volt} (non $\mu$V) e in formato $(N_\text{canali},\, N_\text{campioni})$ (trasposta rispetto alla convenzione CSV). La conversione è:

\begin{formulabox}[Conversione unità di misura]
\[
\mathbf{X}_\text{MNE} = \mathbf{X}_\text{CSV}^{\top} \times 10^{-6}
\qquad \left[\text{shape}: (21,\, N_\text{campioni})\right]
\]
\end{formulabox}

Viene poi creato un oggetto \texttt{mne.Info} contenente i metadati:
\begin{itemize}[noitemsep]
    \item nomi dei canali (\texttt{ch\_names}): lista dei 21 elettrodi
    \item frequenza di campionamento: \texttt{sfreq = 256} Hz
    \item tipo di canale: \texttt{eeg}
\end{itemize}

\begin{lstlisting}[caption={Creazione oggetto MNE Raw}]
data_v = eeg_data.T * 1e-6   # (21, N_campioni), in Volt
info = mne.create_info(ch_names=CH_NAMES, sfreq=256, ch_types="eeg")
raw = mne.io.RawArray(data_v, info)
\end{lstlisting}

% ----
\subsection{Step 2 -- Re-referencing: Common Average Reference (CAR)}

\subsubsection{Problema del riferimento}
Il segnale EEG è sempre misurato come \textbf{differenza di potenziale} tra un elettrodo attivo e un elettrodo di riferimento. Il dataset è stato acquisito con riferimento a \textbf{Pz}. Tuttavia, ogni scelta di riferimento introduce un bias: il potenziale di Pz viene sottratto da tutti i canali, potendo contaminare o mascherare attività locali.

\subsubsection{Common Average Reference}
Il \textit{Common Average Reference} (CAR) ri-referenzia ogni canale rispetto alla media istantanea di tutti gli elettrodi:

\begin{formulabox}[Common Average Reference]
\[
x_i^{\text{CAR}}(t) = x_i(t) - \frac{1}{N} \sum_{j=1}^{N} x_j(t)
\qquad \forall i = 1,\ldots,N
\]
dove $N = 21$ è il numero di canali e $x_i(t)$ è il segnale del canale $i$ al tempo $t$.
\end{formulabox}

\subsubsection{Motivazione}
\begin{itemize}
    \item \textbf{Riduzione artefatti comuni:} rumori presenti su tutti gli elettrodi (es.\ vibrazioni, campo elettrico ambientale) vengono cancellati dalla sottrazione della media.
    \item \textbf{Indipendenza dal sito di riferimento:} il CAR è la scelta standard nella letteratura ERP e nella classificazione EEG-based \cite{nunez1994}.
    \item \textbf{Presupposto:} con 21 canali distribuiti uniformemente sul cuoio capelluto, la somma dei potenziali tende a zero (approssimazione sferica). Con pochi elettrodi questo assunto è meno valido, ma rimane la scelta più robusta disponibile.
\end{itemize}

\begin{lstlisting}[caption={Applicazione CAR in MNE}]
raw_ref, _ = mne.set_eeg_reference(raw, ref_channels="average",
                                    copy=True)
\end{lstlisting}

% ----
\subsection{Step 3 -- Filtraggio Bandpass}

\subsubsection{Motivazione neurofisiologica}
Il segnale EEG grezzo contiene componenti indesiderate a frequenze molto basse (deriva DC, movimenti lenti) e molto alte (rumore elettronico, EMG). Le bande neurofisiologicamente rilevanti per il task decisionale sono:

\begin{table}[h!]
\centering
\caption{Bande di frequenza EEG e ruolo nel task decisionale}
\begin{tabular}{lllp{5cm}}
\toprule
\textbf{Banda} & \textbf{Range (Hz)} & \textbf{Elettrodi chiave} & \textbf{Rilevanza IGT} \\
\midrule
Delta ($\delta$) & 1--4    & Frontale, Parietale & Processi decisionali lenti, attenzione sostenuta \\
Theta ($\theta$) & 4--8    & Fz, F3, F4          & Working memory, controllo cognitivo, monitoraggio errori \\
Alpha ($\alpha$) & 8--13   & O1, O2, P3, P4      & Inibizione cognitiva, resting state \\
Beta ($\beta$)   & 13--30  & Frontale, Centrale  & Elaborazione sensoriale, preparazione motoria \\
Gamma ($\gamma$) & 30--70  & Distribuito         & Binding cognitivo, integrazione informazioni \\
\bottomrule
\end{tabular}
\end{table}

La banda theta frontale è particolarmente rilevante: Cui et al.\ (2013) dimostrano che l'attività theta in Fz aumenta significativamente durante la fase pre-decisionale nell'IGT, riflettendo il carico della working memory e il monitoraggio del rischio.

\subsubsection{Specifiche del filtro}

\begin{formulabox}[Parametri filtro bandpass]
\[
f_\text{low} = 0.5\,\text{Hz} \qquad f_\text{high} = 70\,\text{Hz}
\]
\textbf{Tipo:} FIR (Finite Impulse Response), finestra di Hamming\\
\textbf{Implementazione:} MNE \texttt{raw.filter(l\_freq=0.5, h\_freq=70)}
\end{formulabox}

\paragraph{Perché FIR e non IIR?}
Il paper originale usa un filtro IIR (Butterworth) di 6° ordine. La scelta del filtro FIR nella nostra implementazione è motivata da:
\begin{itemize}
    \item \textbf{Fase lineare:} i filtri FIR con simmetria hanno risposta di fase lineare, preservando la forma temporale dei transienti ERP -- fondamentale per l'analisi delle componenti evento-correlate.
    \item \textbf{Stabilità garantita:} i filtri FIR sono intrinsecamente stabili, a differenza degli IIR.
    \item \textbf{Standard de facto MNE:} la documentazione MNE raccomanda FIR per preprocessing EEG.
\end{itemize}

\paragraph{Taglio basso a 0.5 Hz:} rimuove la deriva DC e le fluttuazioni lente (respirazione $\approx 0.2$ Hz) mantenendo le componenti delta rilevanti.

\paragraph{Taglio alto a 70 Hz:} il paper usa 70 Hz come limite superiore, preservando tutta l'attività corticale rilevante. L'attività muscolare (EMG) è predominante sopra 80--100 Hz.

\begin{lstlisting}[caption={Filtro bandpass FIR}]
raw.filter(l_freq=0.5, h_freq=70.0,
           method="fir", fir_window="hamming")
\end{lstlisting}

% ----
\subsection{Step 4 -- Filtraggio Notch}

\subsubsection{Motivazione}
La rete elettrica introduce un'interferenza sinusoidale alla frequenza di 50 Hz (standard europeo) o 60 Hz (standard nordamericano/messicano). Questa componente non ha origine neurale e deve essere rimossa.

\begin{notabox}[Nota sul dataset]
Il dataset è stato acquisito in Messico (Universidad de Guadalajara), dove la rete elettrica opera a \textbf{60 Hz}. Tuttavia, avendo impostato il notch a \textbf{50 Hz} (standard europeo), è opportuno verificare il segnale PSD del dataset reale. Se si osserva un picco residuo a 60 Hz dopo il preprocessing, il parametro \texttt{NOTCH\_FREQ} va modificato da 50 a 60 nel file \texttt{eeg\_preprocessing.py}.
\end{notabox}

\begin{lstlisting}[caption={Filtro notch}]
raw.notch_filter(freqs=50.0, method="fir")
# Per dataset messicano: sostituire con freqs=60.0
\end{lstlisting}

% ----
\subsection{Step 5 -- Rimozione Artefatti con ICA}

\subsubsection{Il problema degli artefatti EEG}
Il segnale EEG registrato al cuoio capelluto è una sovrapposizione di:
\begin{itemize}
    \item \textbf{Attività neurale:} la componente di interesse
    \item \textbf{Artefatti oculari (EOG):} blink ($\approx 100$--$200\,\mu$V) e movimenti oculari, dominanti sui canali frontali Fp1, Fp2, Fpz
    \item \textbf{Artefatti muscolari (EMG):} contrazioni muscolari, predominanti ad alte frequenze ($>30$ Hz)
    \item \textbf{Artefatti cardiaci (ECG):} battito cardiaco, visibile come picco periodico
\end{itemize}

L'ampiezza degli artefatti da blink (100--200 $\mu$V) è spesso \textbf{10--20 volte} superiore all'ampiezza del segnale EEG corticale (5--20 $\mu$V), rendendoli la principale fonte di inquinamento del segnale.

\subsubsection{Principio dell'ICA}
L'Independent Component Analysis (ICA) è un metodo di separazione cieca delle sorgenti (\textit{Blind Source Separation}). Il modello assunto è:

\begin{formulabox}[Modello ICA]
\[
\mathbf{X} = \mathbf{A}\,\mathbf{S}
\]
dove:
\begin{itemize}[noitemsep]
    \item $\mathbf{X} \in \mathbb{R}^{N_\text{ch} \times T}$ = segnale EEG osservato (21 canali, $T$ campioni)
    \item $\mathbf{S} \in \mathbb{R}^{K \times T}$ = sorgenti indipendenti latenti ($K$ componenti)
    \item $\mathbf{A} \in \mathbb{R}^{N_\text{ch} \times K}$ = matrice di miscelazione (mixing matrix)
\end{itemize}
L'ICA stima $\mathbf{W} = \mathbf{A}^{-1}$ (unmixing matrix) massimizzando la non-gaussianità delle componenti.
\end{formulabox}

La ricostruzione del segnale pulito avviene azzerando le righe di $\mathbf{S}$ corrispondenti agli artefatti:
\[
\hat{\mathbf{X}} = \mathbf{A}_\text{rid}\,\mathbf{S}_\text{rid}
\]
dove il pedice ``rid'' indica la versione ridotta senza le componenti artefattuali.

\subsubsection{Algoritmo: FastICA}
Viene utilizzato l'algoritmo FastICA, che massimizza la non-gaussianità tramite la funzione di contrasto:
\[
G(u) = \log\cosh(u)
\]
I parametri utilizzati sono:
\begin{itemize}[noitemsep]
    \item \textbf{Componenti:} $K = 15$ (con 21 canali, 15 componenti catturano la maggior parte della varianza)
    \item \textbf{Iterazioni massime:} 1500
    \item \textbf{Tolleranza di convergenza:} $10^{-4}$
    \item \textbf{Random state:} 42 (riproducibilità)
\end{itemize}

\subsubsection{Identificazione automatica componenti oculari}
Le componenti artefattuali vengono identificate automaticamente tramite correlazione con i canali EOG frontali (Fp1, Fp2, Fpz):

\begin{lstlisting}[caption={Identificazione e rimozione componenti EOG}]
eog_channels = ["Fp1", "Fp2", "Fpz"]
eog_indices, eog_scores = ica.find_bads_eog(
    raw, ch_name=eog_channels, threshold=3.0
)
ica.exclude = eog_indices
ica.apply(raw_clean)
\end{lstlisting}

Il parametro \texttt{threshold=3.0} indica che vengono escluse le componenti la cui correlazione con il canale EOG supera 3 deviazioni standard rispetto alla distribuzione delle correlazioni. Per il soggetto s-01, le componenti 0 e 1 sono state identificate come artefatti oculari e rimosse.

\begin{notabox}[Nota implementativa -- Pre-filtraggio per ICA]
L'ICA viene fittata su una copia del segnale pre-filtrata con taglio basso a 1 Hz (invece di 0.5 Hz). Questo è raccomandato dalla documentazione MNE: la deriva lenta sotto 1 Hz può compromettere la stima delle componenti indipendenti.
\end{notabox}

% ----
\subsection{Step 6 -- Normalizzazione Z-Score}

\subsubsection{Motivazione}
L'ampiezza del segnale EEG varia considerevolmente tra soggetti diversi (variabilità inter-soggetto). Le cause principali sono:
\begin{itemize}
    \item Diverso spessore del cuoio capelluto e del cranio
    \item Diversa impedenza degli elettrodi
    \item Differenze individuali nell'attività cerebrale di base
\end{itemize}

Senza normalizzazione, un classificatore addestrato su più soggetti tende a imparare le differenze di ampiezza individuali piuttosto che i pattern decisionali.

\subsubsection{Formula}

\begin{formulabox}[Normalizzazione Z-Score per canale]
\[
z_i(t) = \frac{x_i(t) - \mu_i}{\sigma_i}
\]
dove, per ogni canale $i$:
\[
\mu_i = \frac{1}{T}\sum_{t=1}^{T} x_i(t) \qquad
\sigma_i = \sqrt{\frac{1}{T}\sum_{t=1}^{T}(x_i(t) - \mu_i)^2}
\]
Il risultato ha media zero e deviazione standard unitaria \textbf{per ogni canale}.
\end{formulabox}

La normalizzazione è applicata sull'intera registrazione (non sulle singole epoche) per preservare la struttura temporale relativa del segnale.

\begin{lstlisting}[caption={Normalizzazione z-score per canale}]
data = raw.get_data()              # (21, N_campioni)
mu    = data.mean(axis=1, keepdims=True)
sigma = data.std(axis=1, keepdims=True)
sigma[sigma < 1e-10] = 1e-10      # evita divisione per zero
raw._data = (data - mu) / sigma
\end{lstlisting}

% ============================================================
\section{Sincronizzazione EEG--IGT}
% ============================================================

\subsection{Meccanismo di sincronizzazione}

Il file IGT.csv contiene la colonna \texttt{EEG sample} che specifica l'indice esatto del campione EEG corrispondente al momento in cui il partecipante ha effettuato la scelta del deck. Questa colonna è stata registrata in tempo reale durante l'esperimento, garantendo una sincronizzazione precisa a livello di campione.

La conversione in unità di tempo è:

\begin{formulabox}[Conversione campione -- tempo]
\[
t_\text{decisione}^{(k)} = \frac{\text{EEG\_SAMPLE}^{(k)}}{f_s} = \frac{\text{EEG\_SAMPLE}^{(k)}}{256}\quad[\text{secondi}]
\]
dove $k = 1,\ldots,200$ indica il trial.
\end{formulabox}

\subsection{Validazione dei trial}

Un trial è considerato \textbf{valido} per l'estrazione dell'epoca se l'intera finestra temporale $[-2\text{s},\, 0\text{s}]$ ricade all'interno della registrazione EEG disponibile:

\[
\text{VALID}^{(k)} = \begin{cases}
\text{True} & \text{se } \text{EEG\_SAMPLE}^{(k)} \geq 2 \cdot f_s = 512 \\
\text{False} & \text{altrimenti}
\end{cases}
\]

Nel dataset reale tutti i 200 trial per soggetto risultano validi, poiché la baseline iniziale di 3 minuti garantisce un ampio margine prima del primo trial IGT.

\subsection{Generazione delle label}

Le label binarie vengono assegnate automaticamente in base al deck scelto:

\begin{formulabox}[Mapping deck -- label]
\[
\text{label}^{(k)} = \begin{cases}
0 & \text{se } \text{deck}^{(k)} \in \{A, B\} \quad\text{(disadvantageous)}\\
1 & \text{se } \text{deck}^{(k)} \in \{C, D\} \quad\text{(advantageous)}
\end{cases}
\]
\end{formulabox}

\textbf{Motivazione neuroscientifica:} la distinzione A/B vs C/D riflette la dicotomia fondamentale dell'IGT: scelte guidate dalla ricompensa immediata (A/B) vs scelte guidate dall'ottimizzazione a lungo termine (C/D). Questa distinzione è supportata da differenze nell'attivazione prefrontale documentate tramite EEG, fMRI e skin conductance response.

% ============================================================
\section{Estrazione delle Epoche}
% ============================================================

\subsection{Finestra temporale}

\begin{formulabox}[Definizione epoca]
Per ogni trial $k$, l'epoca EEG è definita come:
\[
\mathcal{E}^{(k)} = \mathbf{X}\left[\text{EEG\_SAMPLE}^{(k)} - 2f_s \;:\; \text{EEG\_SAMPLE}^{(k)}\right]
\]
\[
\text{shape: } (21\text{ canali},\; 2\text{s} \times 256\text{ Hz}) = (21,\, 512)
\]
In MNE, la finestra viene definita come $[\texttt{tmin}=-2\text{s},\; \texttt{tmax}=0\text{s}]$.
\end{formulabox}

MNE include il campione corrispondente a $t=0$ nella finestra, producendo $513$ campioni effettivi ($512 + 1$) per ogni epoca.

\subsection{Motivazione della finestra pre-decisionale}

La scelta della finestra $[-2\text{s},\, 0\text{s}]$ è motivata da evidenze neuroscientifiche sulla temporalità dei processi decisionali:

\begin{itemize}
    \item \textbf{Decision Preceding Negativity (DPN):} Bianchin e Angrilli (2011) identificano una negatività frontale in Fz che inizia circa $500$ ms prima della decisione nell'IGT, riflettendo il processo di valutazione del rischio.
    \item \textbf{Slow Cortical Potentials (SCP):} I potenziali corticali lenti si sviluppano su scale temporali di 1--2 secondi, richiedendo una finestra ampia per essere catturati.
    \item \textbf{Attività theta frontale:} L'oscillazione theta (4--8 Hz) in Fz è massima durante la deliberazione, tipicamente nei 1.5--2 secondi che precedono la scelta.
    \item \textbf{Compatibilità con il paradigma:} La finestra non include il feedback (che arriva dopo $t=0$), evitando la contaminazione con componenti post-decisonali come FRN e P300.
\end{itemize}

\subsection{Output delle epoche}

Il risultato finale per ogni soggetto è un tensore tridimensionale:

\begin{formulabox}[Shape del tensore epoche]
\[
\mathbf{X}_\text{subj} \in \mathbb{R}^{N_\text{trial} \times N_\text{ch} \times N_t} = \mathbb{R}^{200 \times 21 \times 513}
\]
\[
\mathbf{y}_\text{subj} \in \{0,1\}^{N_\text{trial}} = \{0,1\}^{200}
\]
\end{formulabox}

% ============================================================
\section{Output e Formato per Machine Learning}
% ============================================================

\subsection{File salvati per soggetto}

Per ogni soggetto vengono salvati quattro file:

\begin{table}[h!]
\centering
\caption{File di output per soggetto}
\begin{tabular}{lll}
\toprule
\textbf{File} & \textbf{Shape} & \textbf{Contenuto} \\
\midrule
\texttt{s-XX\_epochs.npy} & $(200, 21, 513)$ & Segnale EEG normalizzato per ogni trial \\
\texttt{s-XX\_labels.npy} & $(200,)$ & Label binarie (0=disadv, 1=adv) \\
\texttt{s-XX\_samples.npy} & $(200,)$ & Indici campione EEG originali \\
\texttt{s-XX\_info.npz} & -- & Metadati: ch\_names, sfreq, tmin, tmax \\
\bottomrule
\end{tabular}
\end{table}

\subsection{Caricamento per classificazione}

\begin{lstlisting}[caption={Caricamento del dataset completo per ML}]
import numpy as np
from pathlib import Path

epochs_dir = Path('./output/epochs')

all_X, all_y, all_subjects = [], [], []
for label_file in sorted(epochs_dir.glob('*_labels.npy')):
    sid = label_file.stem.replace('_labels', '')
    X_s = np.load(epochs_dir / f'{sid}_epochs.npy')  # (200, 21, 513)
    y_s = np.load(label_file)                          # (200,)
    all_X.append(X_s)
    all_y.append(y_s)
    all_subjects.extend([sid] * len(y_s))

X = np.concatenate(all_X, axis=0)   # (N_tot, 21, 513)
y = np.concatenate(all_y, axis=0)   # (N_tot,)
subjects = np.array(all_subjects)   # per LOSO cross-validation
\end{lstlisting}

\subsection{Strategia di validazione raccomandata}

\begin{notabox}[Leave-One-Subject-Out (LOSO)]
I dati fisiologici sono \textbf{strettamente soggetto-dipendenti}: ogni individuo ha pattern EEG unici legati alla morfologia del cranio, all'attività basale e agli stili cognitivi. Una divisione random train/test che mescola trial dello stesso soggetto produce una stima \textbf{ottimisticamente distorta} delle performance (data leakage inter-soggetto).\\[6pt]
La strategia corretta è la \textbf{Leave-One-Subject-Out (LOSO) cross-validation}: ad ogni fold, il modello viene addestrato su $N-1$ soggetti e testato sul soggetto escluso. Questo stima la capacità di \textbf{generalizzazione a nuovi individui}.
\end{notabox}

\begin{lstlisting}[caption={LOSO cross-validation con scikit-learn}]
from sklearn.model_selection import LeaveOneGroupOut

logo = LeaveOneGroupOut()
for train_idx, test_idx in logo.split(X, y, groups=subjects):
    X_train, X_test = X[train_idx], X[test_idx]
    y_train, y_test = y[train_idx], y[test_idx]
    # ... addestra classificatore su X_train, y_train
    # ... valuta su X_test, y_test
\end{lstlisting}

% ============================================================
\section{Risultati Ottenuti e Confronto con la Letteratura}
% ============================================================

\subsection{Statistiche del dataset preprocessato}

La pipeline è stata eseguita su 55 soggetti (su 59 totali; 4 cartelle risultano mancanti nel dataset scaricato). I risultati sono riassunti nella tabella seguente:

\begin{table}[h!]
\centering
\caption{Statistiche aggregate del dataset preprocessato}
\begin{tabular}{ll}
\toprule
\textbf{Parametro} & \textbf{Valore} \\
\midrule
Soggetti elaborati & 55/59 \\
Trial totali & $55 \times 200 = 11.000$ \\
Epoche per soggetto & 200 (100\% dei trial, tutti validi) \\
Shape tensore per soggetto & $(200, 21, 513)$ \\
Trial advantageous (media) & $\approx 101$ (50.5\%) \\
Trial disadvantageous (media) & $\approx 99$ (49.5\%) \\
Durata media registrazione & $\approx 1056$ s ($\approx 17.6$ min) \\
Componenti ICA rimosse (s-01) & 2 (componenti 0 e 1) \\
\bottomrule
\end{tabular}
\end{table}

\subsection{Distribuzione delle scelte e confronto con il paper}

Il paper (Fig.\ 7) mostra che i partecipanti tendono inizialmente a scegliere deck svantaggiosi (A/B) e progressivamente imparano a preferire quelli vantaggiosi (C/D), in accordo con il modello di apprendimento somatic marker di Bechara.

I dati ottenuti mostrano un'ampia variabilità inter-soggetto nella distribuzione delle scelte:

\begin{itemize}
    \item \textbf{Apprendimento rapido:} soggetti s-51 (167/33), s-27 (159/41), s-17 (160/40) -- scelta prevalentemente vantaggiosa.
    \item \textbf{Apprendimento lento o assente:} soggetti s-36 (59/141), s-13 (62/138), s-03 (66/134) -- scelta prevalentemente svantaggiosa.
    \item \textbf{Distribuzione bilanciata:} maggioranza dei soggetti con proporzioni prossime a 50/50.
\end{itemize}

Questa variabilità è \textbf{pienamente in linea} con la letteratura IGT: Steingroever et al.\ (2013) riportano che una quota significativa di partecipanti sani non sviluppa la preferenza per C/D entro 200 trial.

\subsection{Confronto con la validazione del paper originale}

Il paper effettua una validazione del segnale EEG tramite Power Spectral Density (PSD) e analisi ERP:

\begin{table}[h!]
\centering
\caption{Confronto validazione EEG: paper vs pipeline implementata}
\begin{tabular}{p{4cm}p{4.5cm}p{4.5cm}}
\toprule
\textbf{Aspetto} & \textbf{Paper originale} & \textbf{Pipeline implementata} \\
\midrule
Filtraggio & IIR 6° ordine, 0.5--70 Hz, notch 60 Hz & FIR Hamming, 0.5--70 Hz, notch 50 Hz \\
Riferimento & Auricolari (A1+A2)/2 per ERP & CAR (Common Average Reference) \\
Rimozione artefatti & Non applicata nella validazione & ICA FastICA, 15 componenti \\
Normalizzazione & Non applicata & Z-score per canale \\
Finestra analisi & 200 finestre da 4s (±2s) & [-2s, 0s] pre-decisione \\
Componenti ERP & FRN, P300, N400 identificate & Non analizzate (fase futura) \\
\bottomrule
\end{tabular}
\end{table}

\begin{notabox}[Nota sul notch filter]
Il dataset è stato acquisito in Messico, dove la rete elettrica opera a \textbf{60 Hz}. La pipeline implementata usa 50 Hz (standard europeo). Si raccomanda di verificare la PSD del segnale preprocessato e, se presente un picco residuo a 60 Hz, modificare \texttt{NOTCH\_FREQ = 60.0} nel codice.
\end{notabox}

% ============================================================
\section{Scelte Implementative e Motivazioni}
% ============================================================

\subsection{Libreria MNE-Python}

MNE-Python è stata scelta come framework principale per i seguenti motivi:
\begin{itemize}
    \item \textbf{Standard industriale:} libreria di riferimento per l'analisi EEG/MEG in Python, con oltre 1000 citazioni scientifiche.
    \item \textbf{Gestione completa del dato:} fornisce strutture dati (Raw, Epochs, Evoked) che gestiscono automaticamente metadati, canali e timing.
    \item \textbf{Algoritmi validati:} implementazioni di filtri, ICA, ERP testate dalla comunità scientifica.
    \item \textbf{Integrazione con scikit-learn:} gli array NumPy prodotti da MNE sono direttamente compatibili con la pipeline ML.
\end{itemize}

\subsection{Scelta di $K=15$ componenti ICA}

Con 21 canali disponibili, il numero teorico massimo di componenti indipendenti è 21. La scelta di $K=15$ è un compromesso:
\begin{itemize}
    \item \textbf{Cattura la varianza principale:} 15 componenti tipicamente spiegano $>95\%$ della varianza del segnale.
    \item \textbf{Riduce il rischio di overfitting:} con più componenti, l'ICA ha più difficoltà a convergere e tende a frammentare sorgenti genuine.
    \item \textbf{Accelera il calcolo:} la complessità di FastICA è $O(K^2 \cdot T)$.
\end{itemize}

\subsection{Assenza di baseline correction}

La baseline correction (sottrazione della media in una finestra pre-stimolo) è comune nell'analisi ERP ma \textbf{non è applicata} in questa pipeline. La motivazione è:
\begin{itemize}
    \item La finestra $[-2\text{s},\, 0\text{s}]$ è interamente pre-decisionale: non esiste una finestra ``pre-stimolo'' su cui calcolare la baseline.
    \item La normalizzazione z-score sull'intera registrazione svolge una funzione analoga, rimuovendo il livello DC per canale.
    \item Per l'analisi ERP classica (es.\ replicare i risultati del paper), la baseline correction andrebbe applicata sulla finestra post-feedback $[0\text{s},\, +2\text{s}]$.
\end{itemize}

\subsection{Struttura modulare del codice}

Il codice è organizzato in funzioni indipendenti per garantire:
\begin{itemize}
    \item \textbf{Manutenibilità:} ogni step può essere modificato senza impattare gli altri.
    \item \textbf{Testabilità:} ogni funzione può essere testata in isolamento con dati sintetici.
    \item \textbf{Riutilizzabilità:} le funzioni sono importabili da notebook Jupyter per analisi interattive.
    \item \textbf{Integrabilità:} l'output in formato NumPy è direttamente compatibile con scikit-learn, PyTorch, TensorFlow.
\end{itemize}

% ============================================================
\section{Conclusioni e Sviluppi Futuri}
% ============================================================

\subsection{Riepilogo}

La pipeline implementata esegue in modo completamente automatizzato l'intero preprocessing EEG dal segnale grezzo alle epoche pronte per la classificazione. I punti chiave sono:

\begin{enumerate}
    \item \textbf{Filtraggio robusto:} bandpass FIR 0.5--70 Hz + notch preserva tutte le bande neurofisiologicamente rilevanti.
    \item \textbf{Re-referencing CAR:} riduce la dipendenza dal sito di riferimento e gli artefatti comuni.
    \item \textbf{ICA automatica:} rimozione oggettiva degli artefatti oculari senza intervento manuale.
    \item \textbf{Normalizzazione:} z-score per canale riduce la variabilità inter-soggetto.
    \item \textbf{Sincronizzazione precisa:} uso della colonna EEG\_SAMPLE garantisce allineamento campione-per-campione.
    \item \textbf{Output standard:} formato NumPy compatibile con qualsiasi framework ML.
\end{enumerate}

\subsection{Sviluppi futuri (prossime fasi)}

\begin{itemize}
    \item \textbf{Feature extraction:} calcolo di feature nel dominio del tempo (ampiezza ERP), della frequenza (potenza per banda) e tempo-frequenza (wavelet, Hilbert).
    \item \textbf{Classificazione:} confronto tra SVM, Random Forest, LDA e reti neurali (CNN, LSTM) con cross-validazione LOSO.
    \item \textbf{Replicazione risultati paper:} analisi ERP (FRN, P300, N400) per confronto diretto con la Fig.\ 9 del paper originale.
    \item \textbf{Sensor fusion:} analisi del contributo relativo di sottoinsiemi di canali (es.\ solo frontali vs tutti i canali).
\end{itemize}

% ============================================================
% BIBLIOGRAFIA
% ============================================================
\newpage
\begin{thebibliography}{99}

\bibitem{chavez2026}
Chávez-Sánchez M., Torres-Ramos S., Román-Godínez I., Salido-Ruiz R.A.
\textit{Behavioral and electroencephalographic dataset simultaneously acquired during the Iowa Gambling Task}.
Scientific Data, 13:359, 2026.
\href{https://doi.org/10.1038/s41597-026-06662-0}{DOI: 10.1038/s41597-026-06662-0}

\bibitem{bechara1994}
Bechara A., Damasio A.R., Damasio H., Anderson S.W.
\textit{Insensitivity to future consequences following damage to human prefrontal cortex}.
Cognition, 50:7--15, 1994.

\bibitem{cui2013}
Cui J.F., Chen Y.H., Wang Y., Shum D.H., Chan R.C.
\textit{Neural correlates of uncertain decision making: ERP evidence from the Iowa Gambling Task}.
Frontiers in Human Neuroscience, 7:776, 2013.

\bibitem{bianchin2011}
Bianchin M., Angrilli A.
\textit{Decision preceding negativity in the Iowa Gambling Task: an ERP study}.
Brain and Cognition, 75:273--280, 2011.

\bibitem{nunez1994}
Nunez P.L. et al.
\textit{A theoretical and experimental study of high resolution EEG based on surface Laplacians and cortical imaging}.
Electroencephalography and Clinical Neurophysiology, 90:40--57, 1994.

\bibitem{gramfort2013}
Gramfort A. et al.
\textit{MEG and EEG data analysis with MNE-Python}.
Frontiers in Neuroscience, 7:267, 2013.

\bibitem{steingroever2013}
Steingroever H. et al.
\textit{Performance of healthy participants on the Iowa Gambling Task}.
Psychological Assessment, 25:180, 2013.

\bibitem{kobor2015}
Kóbor A. et al.
\textit{Different strategies underlying uncertain decision making: Higher executive performance is associated with enhanced feedback-related negativity}.
Psychophysiology, 52:367--377, 2015.

\end{thebibliography}

\end{document}
